\section*{Chapter \ref{ch:ml:why-financial-machine-learning-is-different} --- Why Financial Machine Learning Is Different}

\begin{solution}{exo:ml:why-financial-machine-learning-is-different:1}
$R^2 = \rho^2 = 0.25\%$. At $c = kc^\star$, $R^2 = \rho^2(2k - k^2)$; for $k = 3$, $-3\rho^2 = -0.75\%$.
\end{solution}

\begin{solution}{exo:ml:why-financial-machine-learning-is-different:2}
$500/(1 + 499\times0.27) = 3.7$; with 0.05, $500/(1 + 24.95) = 19.3$.
\end{solution}

\begin{solution}{exo:ml:why-financial-machine-learning-is-different:3}
$(2\times0.8/0.5)^2 = 10.2$ months; for three standard errors $(3\times0.8/0.5)^2 = 23.0$.
\end{solution}

\begin{solution}{exo:ml:why-financial-machine-learning-is-different:4}
The ceiling is the R-squared of the truth; a model fitted to the same sample can also fit that sample's noise, which
the truth does not, so its in-sample R-squared can exceed the ceiling. The excess (0.73 points) is fitted noise: it is
the variance term of \cref{prop:ml:why:bv} seen from inside the sample, and it will not recur. Out of sample the trees
cannot beat 0.60\% in expectation and score 0.57\%.
\end{solution}

\begin{solution}{exo:ml:why-financial-machine-learning-is-different:5}
$0.045/\sqrt{240} = 0.29\%$: the window's $-0.56\%$ is 1.9 standard errors from the true zero. A constant error $e$ adds
$e^2$ to the mean squared error: $0.0056^2/0.0093 = 0.34$ points of R-squared. The test decade's returns averaged
$+0.29\%$, so the cross term $2\times0.0056\times0.0029/0.0093 = 0.35$ points doubles the cost: about the 0.67 points
between ridge's 0.38\% and $-0.29\%$.
\end{solution}

\begin{solution}{exo:ml:why-financial-machine-learning-is-different:6}
A model of raw returns learns an intercept equal to the training window's average return, which is the market's path,
estimated from 240 draws; that error alone can make R-squared negative (exercise~5). Train on returns minus each month's
cross-sectional mean, or score the cross-section separately: the same characteristics then give 0.38--0.57\%.
\end{solution}

\begin{solution}{exo:ml:why-financial-machine-learning-is-different:7}
\texttt{strong\_signal(snr=0.0101/0.9899)}: Bayes 1.01\%; ridge 0.67\%, boosted trees 0.42\% (5.1\% in sample), network
0.12\%. The Friedman function is nonlinear, so ridge is biased, but with 20\,000 examples and 1\% of signal the variance
of the flexible models costs more than ridge's bias: at low signal-to-noise the rigid model wins until the sample is
large (the stock panel's trees needed eight years of 500 stocks).
\end{solution}

\begin{solution}{exo:ml:why-financial-machine-learning-is-different:8}
Expand $\E(y - cf)^2 = \sigma_y^2 - 2c\,\Cov(y, f) + c^2\sigma_f^2$ with $\Cov = \rho\sigma_y\sigma_f$; the quadratic in
$c$ peaks at $c^\star = \rho\sigma_y/\sigma_f$ with value $\rho^2\sigma_y^2$, and $R^2(2c^\star) = 0$. A forecast that
correlates at 0.92 with the truth correlates with returns at about $0.92\times\sqrt{0.006} = 0.071$; if its scale is
more than twice $c^\star$ (a model fitted to raw returns carrying a spurious intercept, or one whose predictions spread as
if the signal were strong), $R^2 < 0$ although the ranking of stocks is good. The information coefficient, which ignores
scale, is the better measure of the ranking; R-squared also judges the size.
\end{solution}

\begin{solution}{pb:ml:why-financial-machine-learning-is-different:1}
\begin{enumerate}
\item 0.72\% in the training years, 0.60\% in the test years: the ceiling is a sample statistic of the truth against
realised returns, and the noise (market and specific) differs from decade to decade.
\item $\sqrt{0.0060} = 0.077$.
\item The historical mean of a stock's return is so noisy that beating it is easy; zero is the stricter benchmark (Gu,
Kelly and Xiu).
\item 19\% of the months (23 of 120); the trees lose 18\%.
\item Ridge 0.39\% and 0.38\%; trees 1.45\% and 0.57\%; network 0.64\% and 0.28\%.
\item 64\%, 95\% and 47\%.
\item On the strong task noise is small and bias dominates, so flexible models win and score the same in and out of
sample; on the stocks noise dominates and variance decides.
\item The variance term: the trees have fitted training noise.
\item 3.7: the market and industry factors make the stocks' shocks correlate at 0.27 on average.
\item $-0.29\%$, $-0.09\%$, $-0.30\%$: their intercepts equal the training window's average return, $-0.56\%$.
\item Between 48 months (trees 0.29\%, ridge 0.33\%) and 96 (0.47\% against 0.37\%).
\item 4.3 months from zero; 21 months from ridge.
\item Its R-squared falls to an average of 0.01\% in years 6--10 while the ceiling averages 0.83\%; the inputs, ranks,
look the same before and after.
\item \emph{Named result}: ceiling 0.60\%, the trees reach 95\% of it (0.57\%), 4.3 months to tell them from zero and 21
to tell them from ridge.
\item With ten years, the trees (with a live track long enough to confirm them); with two years, ridge (0.29\% against
0.18\%).
\item The realised R-squared or IC with a sequential test, since the errors, not the inputs, revealed the change.
\item As \texttt{decay\_from}: a characteristic's effect shrinking after it is widely traded.
\item The truth is unobserved; bound it from the best published R-squared or IC of comparable work and from the
spread of models that converge on a level.
\item It is a lower bound on the real ceiling; that very different models cluster between 0.33\% and 0.40\% suggests the
ceiling is not far above.
\item Because the signal is a fraction of a per cent of the variance, changes, is shared across assets and is competed
away.
\end{enumerate}
\end{solution}

\begin{solution}{iq:ml:why-financial-machine-learning-is-different:1}
For individual stocks, monthly, measured against zero out of sample, yes: the best published models reach about
0.4\%. Check the benchmark (zero or mean), the period, whether it is the cross-section or the level, and the
correlation it implies ($\sqrt{0.005} = 0.07$).
\iqlookfor{calibrated expectations of R-squared in finance, and the questions that make a number meaningful.}
\end{solution}

\begin{solution}{iq:ml:why-financial-machine-learning-is-different:2}
Expected squared error at $x$ = noise + squared bias + variance. Returns are almost all noise, so variance is expensive
and a modest bias is cheap: regularised, rigid models are the norm and flexible ones need a lot of data.
\iqlookfor{the decomposition stated correctly and applied to a low signal-to-noise problem.}
\end{solution}

\begin{solution}{iq:ml:why-financial-machine-learning-is-different:3}
Far fewer than ten million. For anything common to all stocks (the market, a factor) there are about 5\,000 days; stocks
are correlated (an effective few dozen per day at best); overlapping labels divide again; and for a slow signal the
independent observations are closer to the number of its half-lives in the sample.
\iqlookfor{dates versus rows, cross-sectional correlation, overlap and persistence.}
\end{solution}

\begin{solution}{iq:ml:why-financial-machine-learning-is-different:4}
The historical mean of an individual stock's return is a very noisy estimate, so a forecast that beats it may have done
so only because the benchmark is bad; zero is a fixed, stricter reference.
\iqlookfor{understanding that the benchmark is a model with estimation error.}
\end{solution}

\begin{solution}{iq:ml:why-financial-machine-learning-is-different:5}
\omterm{def:ml:why-financial-machine-learning-is-different:generalisation}{Overfitting} in the backtest (search size, leakage), drift or decay of the signal (crowding, regime), implementation
differences (data timing, universe, features computed differently live), and plain noise (six months is short). Check
the backtest's trial count and deflated statistics, rerun it on the live period's data, compare features live and
offline, and compute how many months the gap needs to be significant.
\iqlookfor{a structured list and a test for each item, including the possibility that nothing is wrong.}
\end{solution}

\begin{solution}{iq:ml:why-financial-machine-learning-is-different:6}
With mean-zero $y$ and forecast $cf$, $R^2(c) = (2c\rho\sigma_y\sigma_f - c^2\sigma_f^2)/\sigma_y^2$, negative when $c > 2c^\star$
although $\rho > 0$. Shrink the forecast (calibrate its scale on validation data, chapter~11's methods, or Book~7's
isotonic calibration), and judge the ranking by the IC.
\iqlookfor{the formula and the difference between ranking skill and calibration.}
\end{solution}
